% !TeX root = positivity1.tex
\section[Positivity 1: positivity in characteristic 0]{Positivity 1: positivity in characteristic $0$}

Plan:
\begin{enumerate}[label=(\arabic*)]
	\item Positivity in characteristic $0$.
	\item Tools for positivity in characteristic $p$.
	\item Positivity in characteristic $p$.
	\item Discussion, applications, and open questions.
\end{enumerate}


\subsection{Introduction to positivity}

Let $\kkk$ be an algebraically closed field, $X$ a projective variety over $\kkk$.
By divisor, we mean a Cartier divisor.

\paragraph{Positivity of divisors.}
Recall that a divisor $D$ is \emph{ample} if for some $m>0$, the line bundle $\calO_X(mD)$ is very ample, i.e., it gives an embedding of $X$ into some projective space.
Recall the following properties of ample divisors.

\begin{theorem}[Vanishing theorems]
	Let $D$ be an ample divisor on $X$.
	\begin{enumerate}
		\item (Serre vanishing) For any coherent sheaf $\calF$ on $X$, $H^i(X, \calF\otimes \cal_X(mD))=0$ for all $i>0$ and $m\gg 0$.
		\item (Fujita vanishing) For any coherent sheaf $\calF$ on $X$, there exists an integer $m_0$ such that for all $m\ge m_0$ and any nef divisor $N$, we have $H^i(X, \calF\otimes \cal_X(mD+N))=0$ for all $i>0$.
	\end{enumerate}
\end{theorem}

\begin{theorem}[Numerical criterion for ampleness]
	Let $D$ be a divisor on $X$.
	The following are equivalent:
	\begin{enumerate}
		\item $D$ is ample.
		\item (Nakai-Moishezon) For any subvariety $Y\subseteq X$, we have $(D^{\dim Y}\cdot Y)>0$.
		\item (Kleiman) For any non-zero pseudo-effective $1$-cycle $C$ on $X$, we have $(D\cdot C)>0$.
	\end{enumerate}
\end{theorem}

By the numerical criterion for ampleness, we can define the notion of nefness: a divisor $D$ is \emph{nef} if $(D\cdot C)\ge 0$ for any non-zero pseudo-effective $1$-cycle $C$ on $X$,
or equivalently, $(D^{\dim Y}\cdot Y)\ge 0$ for any subvariety $Y\subseteq X$.
And we can define the notion of \emph{strictly nef}: a divisor $D$ is strictly nef if $(D\cdot C)>0$ for any irreducible curve $C$ on $X$.
Also, from the original definition of ampleness, we can define the notion of \emph{semi-ample}:
a divisor $D$ is semi-ample if for some $m>0$, the line bundle $\calO_X(mD)$ is globally generated,
i.e., the global sections of $\calO_X(mD)$ generate the sheaf $\calO_X(mD)$ at every point of $X$.

\begin{remark}
	We have the following implications:
	\begin{align*}
		\text{ample}\implies \text{semi-ample}\implies \text{nef}. \\
		\text{ample}\implies \text{strictly nef}\implies \text{nef}.
	\end{align*}
\end{remark}
\begin{remark}
	Above notations can be generalized to $\bbQ$-Cartier divisors.
\end{remark}
\begin{remark}
	Let $D$ be a nef divisor on $X$ and $A$ an ample divisor on $X$.
	Then $D+tA$ is ample for any $t>0$.
	And hence, all nef divisors are limits of ample divisors.
\end{remark}

\paragraph{Positivity of higher rank sheaves.}
Now we discuss the positivity of higher rank sheaves, for example, vector bundles.
Let $\calE$ be a vector bundle on $X$.
Recall that we can define the projective bundle $\bbP(\calE)$ over $X$ and the tautological line bundle $\calO_{\bbP(\calE)}(1)$ on $\bbP(\calE)$
(here we adopt the notation of Grothendieck, i.e., $\bbP(\calE)$ parameter the quotients of $\calE$).
Then we can define the positivity of $\calE$ by the positivity of $\calO_{\bbP(\calE)}(1)$.
We say that $\calE$ is ample (resp. nef) if $\calO_{\bbP(\calE)}(1)$ is ample (resp. nef).

\begin{definition}[Viehweg's weak positivity]
	Let $\calF$ be a torsion-free coherent sheaf on $X$.
	We say that $\calF$ is \emph{weakly positive} if for any ample divisor $A$ on $X$ and any integer $m>0$, there exists an integer $l>0$ such that $\left(\Sym^{l}(\Sym^{m}(\calF)\otimes \calO_X(A))\right)^{\vee\vee}$ is globally generated on some non-empty open subset of $X$.
\end{definition}

\begin{proposition}
	Let $\calE$ be a vector bundle on $X$.
	Then
	\begin{enumerate}
		\item suppose $\calE$ is nef, then $\calE$ is weakly positive;
		\item if $X$ is a smooth curve, then $\calE$ is nef if and only if $\calE$ is weakly positive.
	\end{enumerate}
\end{proposition}
\begin{proof}
	\textbf{(a)}
	Let $\pi:\bbP(\calE)\to X$ be the projective bundle associated to $\calE$ and $\calO_{\bbP(\calE)}(1)$ the tautological line bundle on $\bbP(\calE)$.
	Then we have $\pi_*\calO_{\bbP(\calE)}(m)=\Sym^m(\calE)$ for any $m>0$.
	Fix ample line bundle $\calH$ on $X$ and integer $m>0$.
	Then $\Sym^m(\calE)\otimes \calH = \pi_*\calO_{\bbP(\calE)}(m)\otimes \calH = \pi_*(\calO_{\bbP(\calE)}(m)\otimes \pi^*\calH)$.
	Note that $\calO_{\bbP(\calE)}(m)\otimes \pi^*\calH$ is ample since $\calO_{\bbP(\calE)}(1)$ is nef and relatively ample.
	Hence $\pi_*(\calO_{\bbP(\calE)}(lm)\otimes \pi^*\calH^{\otimes l}) = \Sym^{l}(\Sym^{m}(\calE)\otimes \calH)$ is globally generated for all sufficiently large $l$ by \cref{rmk:pushforward-of-ample-is-globally-generated}.

	\textbf{(b)}
	Let $C$ be a smooth curve and $\calE$ a vector bundle on $C$.
	Suppose $\calE$ is weakly positive.
	Let $D \subset \bbP(\calE)$ be an irreducible curve surjecting onto $C$ via $\pi:\bbP(\calE)\to C$.
	Denote by $f: D \to C$ the restriction of $\pi$ to $D$.
	Then we have


	For (b),


	\Yang{}.
\end{proof}

\begin{remark}\label{rmk:pushforward-of-ample-is-globally-generated}
	Let $f: X \to Y$ be a proper fibration of projective varieties over $\kkk$ and $\calL$ is an ample line bundle on $X$.
	Then $f_*\calL^{\otimes n}$ is globally generated for sufficiently large $n$.

	First we show that there exists an integer $n_0$ such that restricted map $H^0(X, \calL^{\otimes n}) \to H^0(X_y, \calL^{\otimes n}|_{X_y})$ is surjective for any $n\ge n_0$ and any $y\in Y$.
	Consider the exact sequence
	\[ 0 \to \calI_{X_y}\otimes \calL^{n} \to \calL^{\otimes n} \to \calL^{\otimes n}|_{X_y} \to 0. \]
	Then we have the long exact sequence
	\[ H^0(X, \calI_{X_y}\otimes \calL^{\otimes n}) \to H^0(X, \calL^{\otimes n}) \to H^0(X_y, \calL^{\otimes n}|_{X_y}) \to H^1(X, \calI_{X_y}\otimes \calL^{\otimes n}). \]
	By Serre vanishing, there exists $n_y$ such that $H^1(X, \calI_{X_y}\otimes \calL^{\otimes n})=0$ for any $n\ge n_y$.
	By cohomology and proper base change theorem, there is an open neighborhood $U_y$ of $y$ such that $H^1(X, \calI_{X_y}\otimes \calL^{\otimes n})=0$ for any $n\ge n_y$ and any $y'\in U_y$.
	By the quasi-compactness of $Y$, there is $n_0$ such that $H^1(X, \calI_{X_y}\otimes \calL^{\otimes n})=0$ for any $n\ge n_0$ and any $y\in Y$.
	Then the map $H^0(X, \calL^{\otimes n}) \to H^0(X_y, \calL^{\otimes n}|_{X_y})$ is surjective for any $n\ge n_0$.

	Then we show that $f_*\calL^{\otimes n}$ is globally generated for any $n\ge n_0$.
	Let $y \in Y$ be a closed point.
	Recall that the natural map $(f_*\calL^{\otimes n})\otimes \rkk(y) = f_*(\calL^{\otimes n})|_{\{y\}} \to H^0(X_y, \calL^{\otimes n}|_{X_y})$ is given by the restriction map of global sections to the fiber $X_y$.
	Hence the map $H^0(X, \calL^{\otimes n}) \to H^0(X_y, \calL^{\otimes n}|_{X_y})$ factors through the map $H^0(Y, f_*\calL^{\otimes n}) \to f_*(\calL^{\otimes n})$.
	Then by the surjectivity of $H^0(X, \calL^{\otimes n}) \to H^0(X_y, \calL^{\otimes n}|_{X_y})$ and the isomorphism $H^0(X, \calL^{\otimes n}) \cong H^0(Y, f_*\calL^{\otimes n})$, we have the surjectivity of $H^0(Y, f_*\calL^{\otimes n}) \to f_*(\calL^{\otimes n})\otimes \rkk(y)$ for any $n\ge n_0$.
	By the Nakayama lemma, the map $H^0(Y, f_*\calL^{\otimes n})\otimes \calO_{Y,y} \to f_*\calL^{\otimes n}\otimes \calO_{Y,y}$ of stalks is surjective for any $n\ge n_0$ and any closed point $y\in Y$.
	Note that the surjective locus of the map $H^0(Y, f_*\calL^{\otimes n})\otimes \calO_{Y} \to f_*\calL^{\otimes n}$ is open since it is the complement of the support of the cokernel of the map.
	Then by the density of closed points, we have the surjectivity of $H^0(Y, f_*\calL^{\otimes n})\otimes \calO_{Y} \to f_*\calL^{\otimes n}$ for any $n\ge n_0$.
\end{remark}
\begin{remark}
	The \cref{rmk:pushforward-of-ample-is-globally-generated} does not hold if we replace the ample line bundle $\calL$ by a globally generated line bundle $\calL$.

	Let $f: X \to \bbP^1$ be a non-trivial elliptic fibration with zero section $S$.

	\Yang{}

\end{remark}





\subsection{Canonical divisor and duality}

\begin{theorem}[{Grothendieck duality, \cite[Tags \href{https://stacks.math.columbia.edu/tag/0A9Y}{0A9Y} and \href{https://stacks.math.columbia.edu/tag/0AU3}{0AU3}]{Stacks}}]
	Let $f:X\to Y$ be a proper morphism of schemes of finite type over a noetherian base schemes.
	Then there exists a functor $f^!:\sfD^+(Y)\to \sfD^+(X)$ such that for any $\calF\in \sfD^-(X)$ and $\calG\in \sfD^+(Y)$, we have
	\[ \sfR f_*\sfR \calHom_X(\calF, f^!\calG) \cong \sfR \calHom_Y(\sfR f_*\calF, \calG). \]
\end{theorem}

\begin{remark}
	Recall some fact about the derived category (\cite[Tag \href{https://stacks.math.columbia.edu/tag/05QI}{05QI}]{Stacks}).
	For any abelian category $\bfA$, we can define the derived category $\sfD(\bfA)$ of $\bfA$.
	The objects of $\sfD(\bfA)$ are the complexes of objects in $\bfA$ and the morphisms are extanded by localization by quasi-isomorphic, i.e., if a morphism $f: A^\bullet \to B^\bullet$ induces isomorphisms on cohomology, then $f$ is an isomorphism in $\sfD(\bfA)$.
	In particular, if $\bfA$ has enough injectives (resp. projectives), then an object $A \in \Obj(\bfA) \subset \Obj(\sfD(\bfA))$ is isomorphic to an injective (resp. projective) resolution of it in $\sfD(\bfA)$.
	There are some subcategories of $\sfD(\bfA)$, namely, $\sfD^+(\bfA)$, $\sfD^-(\bfA)$, and $\sfD^b(\bfA)$, which are the full subcategories of $\sfD(\bfA)$ consisting of complexes has bounded below, bounded above, and bounded cohomology, respectively.
	Let $X$ be a scheme and $\bfA = \QCoh(X)$ be the category of quasi-coherent sheaves on $X$.
	We denote by $\sfD(X) = \sfD(\QCoh(X))$ (and similarly for $\sfD^+(X)$, $\sfD^-(X)$, and $\sfD^b(X)$) here.

	Let $\bfA$ and $\bfB$ be abelian categories and $F: \bfA \to \bfB$ an additive functor.
	Suppose that $\bfA$ has enough injectives and $F$ is left exact.
	Then we can define its \emph{right derived functor} $\sfR F: \sfD(\bfA) \to \sfD(\bfB)$ of $F$.
	For an complex $I^\bullet = (\cdots \to I^n \to I^{n+1} \to \cdots)$ consisting of injective objects in $\bfA$, we have $\sfR F(I^\bullet) = (\cdots \to F(I^n) \to F(I^{n+1}) \to \cdots)$.
	In particular, we can compute $\sfR F(A)$ for an object $A \in \Obj(\bfA)$ by taking an injective resolution.
	For $\calF, \calG \in \sfD(X)$, the functor $\sfR \calHom_X(\calF, \calG)$ above can be defined as the right derived functor of the left exact functor $\calHom_X(\calF, -): \QCoh(X) \to \Mod_{\calO_X}$ (or $\calHom_X(-, \calG): \Coh(X) \to \QCoh(X)$).\Yang{ref required}
\end{remark}

\begin{example}
	Suppose that $f: X \to Y$ is a smooth projective morphism of relative dimension $n$.
	Then we have $f^!\calO_Y \cong \omega_{X/Y}[n]$ with $\omega_{X/Y} = \bigwedge^n \Omega_{X/Y}$ the relative canonical sheaf.
	Here $A^\bullet[n]$ is the complex with $A^i$ in degree $i-n$ for a complex $A^\bullet$.
	Hence $\omega_{X/Y}[n]$ is a complex with $\omega_{X/Y}$ in degree $-n$ and $0$ in other degrees.
	\Yang{}
\end{example}

Many other duality theorems can be deduced from Grothendieck duality.
Let us consider the Serre duality as an example.

\begin{example}
	Let $\calF$ be a locally free sheaf on a smooth projective variety $X$ of dimension $n$ over an algebraically closed field $\kkk$.
	Take $\calG = \calO_Y = \kkk$.
	Then the Grothendieck duality gives
	\[ \sfR \Gamma(\sfR \calHom_X(\calF, f^!\calO_Y)) \cong \sfR \Hom_\kkk(\sfR \Gamma \calF, \kkk). \]
	\[ \Hom_\kkk(\sfR \Gamma(\calF), \kkk) \cong \sfR \Gamma(\sfR \calHom_X(\calF, f^!\calO_Y)). \]
	The sheaf $f^!\calO_Y$ is just the twisted cannonical sheaf $\omega_X[n]$ with $\omega_X = \bigwedge^n \Omega_X$ since $X$ is projective smooth.
	The LHS is equal to $\Hom_\kkk(\sfR \Gamma(\calF), \kkk) = \Hom(H^i(X, \calF), \kkk)[-i]$.\Yang{}.
	Then RHS is equal to $\sfR \Gamma(\sfR \calHom_X(\calF, \omega_X[n])) = \sfR \Gamma(\calF^\vee\otimes \omega_X[n]) = H^{n-i}(X, \calF^\vee\otimes \omega_X)[-i]$.\Yang{}
	\Yang{to be continued}
\end{example}

Set the dualing complex $\omega_{X/Y} = f^!\calO_Y[-n]$.

For general $\calG$, $f^! \calG \cong \calG \otimes^\sfL f^!\calO_Y$.


Let $f: X \to Y$ be a finite surjective morphism of projective varieties over $\kkk$.






\subsection{Some questions}

Recall that we can define the Iitaka dimension of a divisor $D$ on a projective variety $X$.

\begin{conjecture}[Abundance conjecture]
	Let $X$ be a smooth projective variety over an algebraically closed field $\kkk$ and $D$ a nef divisor on $X$.
	If the Iitaka dimension of $D$ is equal to the numerical dimension of $D$, then $D$ is semi-ample.
	\Yang{}
\end{conjecture}


\begin{conjecture}[Iitaka conjecture]
	Let $f:X\to Y$ be a surjective morphism of smooth projective varieties over $\bbC$.
	Then we have
	\[ \kappa(X) \ge \kappa(Y) + \kappa(F), \]
	where $F$ is a general fiber of $f$.
	\Yang{}

\end{conjecture}

When $Y$ is of general type, Viehweg proved the Iitaka conjecture by using the notion of weak positivity.


\subsection[Positivity results in characteristic 0]{Positivity results in characteristic $0$}

\begin{theorem}[Viehweg' 83]
	Let $f:X\to Y$ be a surjective projective morphism of smooth projective varieties over $\bbC$.
	Then
	\begin{enumerate}
		\item $f_*\omega_{X/Y}$ is weakly positive;
		\item $f_*\omega_{X/Y}^{\otimes m}$ is weakly positive for any $m>0$.
	\end{enumerate}
\end{theorem}

\begin{theorem}[Campana'04]
	Let $(X,\Delta)$ be a lc pair and $f:X\to Y$ a surjective projective morphism of smooth projective varieties over $\bbC$.
	Then $f_*\calO_X(m(K_{X/Y}+\Delta))$ is weakly positive for any $m>0$ such that $m(K_{X/Y}+\Delta)$ is Cartier.
	\Yang{}
\end{theorem}

The key ingredient of the proof is
\begin{enumerate}
	\item VHS (Fujino-Fujisana 18) / Kodaira vanishing + Mumford regularity.
	\item (a) + covering trick.
\end{enumerate}

\begin{example}
	Let $f: X \to C$ be a fibration with $C$ a smooth projective curve of genus $g(C) \ge 2$.
	\Yang{}
\end{example}
























